% TOP_PROBLEM: {"id": 6900001, "title": "Configuration Spaces of Tensegrities — Problem 1", "category_id": 6, "category": {"id": 6, "name": "geometry", "display_name": "Geometry", "description": "Euclidean and non-Euclidean geometry, geometric structures.", "slug": "geometry", "order_index": 6, "created_at": "2026-07-31T15:26:25.671Z"}, "status": "partially_solved"}
% FIRST_POSTED: null
% ENTRY_KIND: statement_only
% Imported from: batch18.tex; UnsolvedMath ID 6900001
\documentclass[11pt]{article}
\usepackage[margin=1in]{geometry}
\usepackage{amsmath,amssymb,mathtools}
\usepackage{fontspec,unicode-math}
\setmainfont{Latin Modern Roman}
\setsansfont{Latin Modern Sans}
\setmathfont{Latin Modern Math}
\usepackage{xurl,hyperref}
\hypersetup{colorlinks=true,urlcolor=blue,linkcolor=blue}
\newcommand{\sref}[2]{\hyperlink{source:#1:#2}{[S#2]}}
\newcommand{\cataloguescope}[1]{\paragraph{Recorded mathematical question.}#1}
\begin{document}
% UnsolvedMath ID 6900001; source catalogue code AMR-068-0001
\setcounter{section}{6900000}
\section{Configuration Spaces of Tensegrities — Problem 1}\label{6900001}
{\small\sffamily UnsolvedMath ID 6900001 · Geometry}
\subsection{Definitions and mathematical statement}

\paragraph{Shared notation.}$\mathbb N=\{1,2,\ldots\}$, and $\mathbb Z,\mathbb Q,\mathbb R,\mathbb C$ denote the integers, rationals, reals and complex numbers. Any other conventions are specified locally.


Let $G$ be a simple graph with labeled vertices $1,\ldots,n$. A configuration is $P=(p_1,\ldots,p_n)\in(\mathbb R^d)^n$; coincident vertices are allowed. A self-stress is a real edge assignment $w_{ij}=w_{ji}$, set to zero on missing edges, satisfying equilibrium at every vertex:
\[\sum_{j\ne i}w_{ij}(p_j-p_i)=0\quad(1\leq i\leq n).\]
Let $W(G,P)$ be the vector space of all such self-stresses. Negative and positive edge values are the source’s cable and strut signs. The base $B_d(G)$ is the full labeled configuration space $(\mathbb R^d)^n$, stratified by the following fiber equivalence: $W(G,P)$ and $W(G,Q)$ are equivalent if a homeomorphism between them preserves the sign of every edge coordinate of every stress. A stratum is a maximal connected set of configurations with equivalent fibers; these strata are semialgebraic. This base is not quotiented by Euclidean or affine motions. Write $K_n$ for the complete graph on $n$ labeled vertices.
\paragraph{Statement recorded in the supplied source.}
Describe the combinatorics of the strata of $B_2(K_6)$, $B_3(K_4)$ and $B_3(K_5)$, in the manner of the explicit stratified descriptions of $B_2(K_4)$ and $B_2(K_5)$ in the source. \sref{6900001}{2}
\subsection{Short English statement}
Describe the stratified configuration spaces for six planar vertices and for four or five spatial vertices, including their stratum combinatorics.
\subsection{Sources}
\begin{description}
\item[\hypertarget{source:6900001:1}{S1}] ULAM.AI and the UnsolvedMath Contributors, \emph{UnsolvedMath: A Curated Collection of Open Mathematics Problems}, record 6900001 (AMR-068-0001). CC BY 4.0. \url{https://huggingface.co/datasets/ulamai/UnsolvedMath}.
\item[\hypertarget{source:6900001:2}{S2}] Oleg Karpenkov, Jan Schepers and Brigitte Servatius, \emph{On stratifications for planar tensegrities with a small number of vertices}, arXiv:1201.3557v1 (2012), Definitions 1.1--1.3, pp. 2--3; Problem 2.7, p. 12. \url{https://arxiv.org/pdf/1201.3557}.
\end{description}
\label{end:6900001}
\end{document}
